\subsection*{把好听的解释变成可以检验的假设}
我们容易把隐藏神经元想象成边缘、笔画或概念检测器，但这种解释必须与实际参数和行为相符。先固定网络，找出什么输入能强烈激活某个单元，再查看看似相同的特征在不同背景下是否仍产生类似响应。一次漂亮的示例不足以建立稳定解释。

把输入稍微改变，观察输出怎样变，是更直接的实验。仿射层满足 $\Delta z=W\Delta x$：偏置在相减时消失，所以敏感性由权重和扰动方向共同决定。某方向的小变化可能被放大，另一些方向可能被忽略。

\subsection*{准确、稳定与可信不是同一个量}
即使原图被正确分类，平移、旋转或加一点扰动后也可能变化；普通全连接层并不会自动知道这些变化应保持类别。高置信度也只是模型输出分布尖锐，并不保证它遇到的是熟悉数据，更不保证解释正确。

设计检查时可以分三步：保留独立测试样本观察泛化；做有控制的扰动检验稳定性；对置信度分组比较预测与实际频率。三步回答不同问题，不应让一个总体准确率代替所有分析。

本节的小线性算例让你能精确追踪一次扰动，但真实网络的解释还涉及非线性与数据分布。读课程中的可视化时，应同时问“这张图支持了什么结论”和“还排除了哪些替代解释”。数学公式与实验证据在这里承担不同角色。
